\documentclass[11pt, oneside]{article}
\usepackage[margin=0.75in]{geometry}
\geometry{letterpaper}
\usepackage{graphicx}
\usepackage{amssymb}
\usepackage{amsmath}
\usepackage{booktabs}
\usepackage{hyperref}

\hypersetup{
    colorlinks=true,
    linkcolor=blue,
    filecolor=magenta,
    urlcolor=cyan,
    citecolor=blue,
}

\setlength{\parskip}{0.1in}
\setlength{\parindent}{0in}

\title{CS779 Project: Neural Machine Translation for English--Bengali and English--Hindi}

\author{
    Kunal Chandra \\
    240580 \\
    {\tt kunalc24@iitk.ac.in} \\
    {Indian Institute of Technology Kanpur}
}

\date{}

\begin{document}
\maketitle

\begin{abstract}
This report documents the development of a Neural Machine Translation (NMT) system for
English-to-Bengali and English-to-Hindi translation. Development proceeded through four
weekly iterations. I began with a vanilla GRU encoder--decoder as a baseline, then built a
preprocessing pipeline using fastText embeddings, spaCy for English, and
\texttt{indic-nlp-library} for the target scripts. From there the architecture moved to a
bidirectional LSTM with Bahdanau attention, then to a vanilla Post-LN Transformer, and then
to a stacked bidirectional GRU with global attention. For the testing phase I returned to the
Transformer family with a Pre-Layer-Normalization design, weight tying, and GELU activations,
which produced the best results. A byte-pair-encoding variant was also built during the
testing phase; it segmented Bengali poorly and scored substantially worse despite a lower
training loss, so the final submission used the word-level model.
\end{abstract}

\section{Competition Result}

\begin{table}[h!]
\centering
\small
\begin{tabular}{@{}ll@{}}
\toprule
Item & Value \\
\midrule
CodaLab username & \texttt{k\_240580} \\
Final rank & 35 \\
BLEU & 0.162 \\
chrF & 0.417 \\
ROUGE-L & 0.445 \\
Training-phase submissions & 25 \\
Testing-phase submissions & 7 \\
\bottomrule
\end{tabular}
\caption{Final competition standing.}
\end{table}

\begin{table}[h!]
\centering
\small
\begin{tabular}{@{}lc@{}}
\toprule
Week & Submissions \\
\midrule
Week 1 & 2 \\
Week 2 & 4 \\
Week 3 & 7 \\
Week 4 & 12 \\
\bottomrule
\end{tabular}
\caption{Training-phase submissions by week.}
\end{table}

\section{Problem Description}

The competition required building an NMT system for English-to-Bengali and English-to-Hindi
translation. Given a training set of sentence pairs and a set of test source sentences, the
objective was to generate translations evaluated using BLEU, chrF++, and ROUGE-L. The
principal challenges were:

\begin{itemize}
    \item extracting and aligning the data from the supplied JSON files;
    \item handling large target vocabularies under GPU memory constraints;
    \item preventing overfitting on a comparatively small parallel corpus;
    \item keeping evaluation metrics reliable and comparable across weeks;
    \item using the available GPU efficiently, since training time was the binding constraint
          on how many architectures could be tried.
\end{itemize}

\section{Data Analysis}

\subsection{Dataset Statistics}

\begin{table}[h!]
\centering
\small
\begin{tabular}{@{}lrrrr@{}}
\toprule
Pair & Train & Validation & Test & Avg.\ length \\
\midrule
English--Bengali & 68{,}849 & 9{,}836 & 19{,}672 & 16.8 \\
English--Hindi & 149{,}646 & 21{,}379 & 42{,}757 & 31.4 \\
\bottomrule
\end{tabular}
\caption{Corpus sizes. Average sentence length was stable across splits (Bengali 16.8 / 16.87
/ 16.9 and Hindi 31.4 / 31.45 / 31.5 for train / validation / test), indicating the splits
were drawn from the same distribution.}
\end{table}

\begin{table}[h!]
\centering
\small
\begin{tabular}{@{}lrr@{}}
\toprule
Pair & Source vocabulary & Target vocabulary \\
\midrule
English--Bengali & 31{,}920 & 37{,}921 \\
English--Hindi & 33{,}366 & 31{,}680 \\
\bottomrule
\end{tabular}
\caption{Vocabulary sizes after preprocessing, with a frequency-2 cutoff.}
\end{table}

The Hindi corpus has roughly twice the average sentence length of the Bengali corpus, which
matters more than the raw sentence count: attention cost grows quadratically in sequence
length, so the Hindi pair was consistently the more expensive of the two to train.

\subsection{Critical Data Issues}

The raw text contained a large quantity of junk tokens, most visibly long digit sequences and
alphanumeric tracking codes. These carry no translatable meaning but do inflate the
vocabulary, so they were removed with regular expressions during preprocessing. Left in, they
would have consumed embedding capacity on tokens the model can never usefully translate.

\section{Model Description}

\subsection{Week 1: Vanilla GRU (Baseline)}

\textbf{Architecture:} a simple encoder--decoder RNN using GRUs, with manual sequential
decoding and teacher forcing during training.

\textbf{Goal:} produce a first submission and establish a minimal working baseline.

\textbf{Key learning:} this is the simplest recurrent setup. It trains quickly but compresses
the entire source sentence into a single fixed-size vector, so it loses context on longer
inputs.

\subsection{Week 2: Bidirectional LSTM with Bahdanau Attention}

\begin{enumerate}
    \item \textbf{Model:} a bidirectional LSTM encoder--decoder with Bahdanau attention.
    \item \textbf{Speed fix --- dynamic padding:} this week was largely spent on latency. I
          introduced \texttt{pack\_padded\_sequence} and \texttt{pad\_packed\_sequence}, which
          ensure the recurrent layers compute hidden states only for non-padding tokens. This
          removed a substantial amount of wasted computation.
    \item \textbf{Architecture upgrade:} the unidirectional GRU became a bidirectional LSTM,
          so the encoder sees both left and right context for every source token.
    \item \textbf{Parameter changes:} embedding dimension 300 and hidden dimension 256. I set
          \texttt{embedding.weight.requires\_grad = False} to freeze the pre-trained vectors,
          on the assumption that the pre-trained knowledge was worth preserving.
    \item \textbf{Embeddings:} 300-dimensional fastText vectors were used as static
          embeddings.
\end{enumerate}

\subsection{Week 3: Vanilla Transformer}

\textbf{Architecture:} the standard \texttt{nn.Transformer} in its Post-Layer-Normalization
form, with 6 encoder and 6 decoder layers, trained under a \texttt{LambdaLR} schedule.

\textbf{Key learnings:} the latency problem was solved outright --- unlike an RNN, the
Transformer processes all positions in parallel, and multi-head attention handles long-range
dependencies far better. Convergence, however, was slow and the Post-LN arrangement proved
sensitive to the learning rate, which motivated the Pre-LN design adopted later.

\subsection{Week 4: Stacked Bidirectional GRU with Attention}

\textbf{Architecture:} a stacked bidirectional GRU encoder paired with a unidirectional GRU
decoder using global attention.

\textbf{Key learnings:} this was the strongest recurrent model of the four weeks, but it
confirmed that the RNN family is limited by its sequential nature and by the attention
bottleneck. Adding global attention helped, yet the architecture could not match a
Transformer's throughput on the same hardware budget.

\subsection{Final Model: Pre-LN Transformer}

The final architecture combines structural and regularization choices aimed at stability
rather than raw capacity.

\begin{itemize}
    \item \textbf{Depth:} 6-layer encoder and 6-layer decoder stacks.
    \item \textbf{Pre-LN:} layer normalization is applied \emph{before} each attention and
          feed-forward block, which keeps gradient magnitudes well behaved at initialization
          and makes the warmup schedule usable [2].
    \item \textbf{Activation:} the Gaussian Error Linear Unit (GELU), standard in modern
          Transformer stacks and smoother than ReLU through a deep network [5].
    \item \textbf{Weight tying:} the target embedding matrix is shared with the final output
          projection [3]. This cut the parameter count appreciably and was the main defense
          against overfitting on a small corpus.
\end{itemize}

The implementation draws on three results. The multi-head attention mechanism and the overall
encoder--decoder design come from \emph{Attention Is All You Need} [1]. The Pre-LN
arrangement follows \emph{On Layer Normalization in the Transformer Architecture} [2], which
was what finally allowed stable training at a high peak learning rate. The weight-tying
technique is taken from \emph{Tying Word Vectors and Word Classifiers} [3].

Decoding is greedy throughout. Beam search would likely have improved the scores, but it was
too slow to fit inside the submission schedule.

During the testing phase I also built a pipeline using a BPE tokenizer. It produced poor
segment boundaries on Bengali: training loss went down while BLEU went down with it, which is
the clearest possible sign that the loss was measuring the wrong thing. That variant was
dropped in favour of the word-level model.

After the final submission I continued training the same architecture for more epochs and it
scored materially better, but the deadline had passed. Those predictions are included in the
accompanying archive.

\begin{figure}[htbp]
     \centering
     \includegraphics[width=\textwidth]{Architecture.png}
     \caption{Final model architecture.}
     \label{fig:model}
\end{figure}

\section{Experiments and Training}

\subsection{Data Preprocessing}

The two sides of the corpus fail in different ways, so the pipeline is deliberately
asymmetric.

\subsubsection{Source Language (English)}

\begin{enumerate}
    \item \textbf{Normalization and contraction expansion:} raw text is stripped of URLs,
          HTML tags, and proprietary encodings. A custom routine expands contractions
          (\texttt{n't} $\rightarrow$ \texttt{ not}, \texttt{'s} $\rightarrow$ \texttt{ is}),
          so that informal and formal spellings collapse to one vocabulary entry.
    \item \textbf{Tokenization:} spaCy's pipeline, with the parser and tagger disabled since
          only token boundaries are needed.
    \item \textbf{Junk filtering:} pure punctuation, currency symbols, and alphanumeric junk
          such as tracking codes are discarded. This prevents vocabulary bloat.
\end{enumerate}

\subsubsection{Target Languages (Bengali and Hindi)}

\begin{enumerate}
    \item \textbf{Noise removal:} as with English, URLs, HTML, and stray punctuation are
          removed.
    \item \textbf{Character normalization:} \texttt{IndicNormalizerFactory} standardizes the
          script. This step matters more than it appears --- Indic scripts admit several
          encodings of the same vowel or conjunct, and without normalization one word appears
          as several distinct vocabulary entries.
    \item \textbf{Tokenization:} \texttt{indic\_tokenize.trivial\_tokenize} for token
          boundaries in the target script.
    \item \textbf{Final filtering:} tokens are restricted to valid Bengali/Hindi characters
          and numerals.
\end{enumerate}

The same pipeline was applied to the training, validation, and test splits, so that no
distribution shift is introduced by preprocessing alone.

\subsection{Training Configuration}

\subsubsection{Week 1: Vanilla GRU}

Embedding size 128, Adam, learning rate $3 \times 10^{-3}$, \texttt{nn.NLLLoss}, batch size
64, greedy decoding, 8 epochs.

\textbf{Justification:} the small embedding minimized memory overhead and the high learning
rate gave rapid initial convergence. The configuration existed only to establish a verifiable
baseline and to expose the latency limits of sequential decoding.

\subsubsection{Week 2: Bidirectional LSTM with Attention}

Embedding size 300, hidden size 256, AdamW, learning rate $3 \times 10^{-3}$,
\texttt{nn.CrossEntropyLoss}, batch size 32, greedy decoding, dropout 0.3, 8 epochs and then
22.

\textbf{Justification:} capacity was increased to make attention worthwhile, and dropout was
introduced because the Week 1 model had shown no regularization at all. Out-of-memory errors
and low GPU utilization were the main reasons for moving on.

\subsubsection{Week 3: Vanilla Transformer}

Model dimension 512, feed-forward dimension 2048, 8 attention heads, 6 encoder and 6 decoder
layers, Adam, cross-entropy loss, \texttt{LambdaLR} with 4000 warmup steps.

\textbf{Justification:} these are the proportions of the original ``base'' Transformer, paired
with the inverse-square-root warmup schedule the architecture was designed around.

\subsubsection{Week 4: Stacked GRU}

Embedding size 256, hidden dimension 512, 2 layers, dropout 0.3, cross-entropy loss.

\textbf{Justification:} this was the last recurrent configuration tried. The published
evidence in favour of Transformers, combined with the throughput gap measured in Week 3,
motivated the return to the Transformer family for the testing phase.

\subsubsection{Testing Phase: Final Model}

\begin{table}[h!]
\centering
\small
\begin{tabular}{@{}ll@{}}
\toprule
Hyperparameter & Value \\
\midrule
Embedding / model dimension & 512 \\
Attention heads & 8 \\
Feed-forward dimension & 1024 \\
Encoder / decoder layers & 6 / 6 \\
Dropout & 0.15 \\
Maximum sequence length & 55 \\
Optimizer & AdamW ($\beta = 0.9/0.98$, weight decay 0.01) \\
Learning-rate schedule & Noam inverse square root, 4000 warmup steps \\
Loss & Cross-entropy, label smoothing 0.1 \\
Precision & Mixed (fp16 autocast with gradient scaling) \\
Decoding & Greedy \\
Epochs & 8 for the submitted model; 15 subsequently \\
\bottomrule
\end{tabular}
\caption{Final model configuration.}
\end{table}

Raising the epoch count from 8 to 15 improved results noticeably, but that run finished after
the submission deadline.

\section{Results}

\subsection{Validation Set (Training Phase)}

\begin{table}[h!]
\centering
\small
\begin{tabular}{@{}lccc@{}}
\toprule
Model & chrF & ROUGE-L & BLEU \\
\midrule
Vanilla GRU (baseline) & 0.233 & 0.319 & 0.058 \\
Bi-LSTM with attention & 0.264 & 0.277 & 0.078 \\
Vanilla Transformer (Post-LN) & 0.280 & 0.340 & 0.073 \\
Stacked GRU with global attention & 0.326 & 0.362 & 0.102 \\
\bottomrule
\end{tabular}
\caption{Validation performance by week.}
\end{table}

The Week 3 Transformer is the interesting row: it beat the Bi-LSTM on chrF and ROUGE-L while
scoring \emph{lower} on BLEU. BLEU rewards exact $n$-gram matches, so a model producing
broadly correct content with different phrasing can move the three metrics in opposite
directions.

\subsection{Test Set}

\begin{table}[h!]
\centering
\small
\begin{tabular}{@{}lccc@{}}
\toprule
Model & chrF & ROUGE-L & BLEU \\
\midrule
Pre-LN Transformer with BPE & 0.192 & 0.218 & 0.059 \\
\textbf{Pre-LN Transformer, word-level (submitted)} & \textbf{0.417} & \textbf{0.445} & \textbf{0.162} \\
\bottomrule
\end{tabular}
\caption{Test performance.}
\end{table}

The word-level Pre-LN Transformer was the strongest model on every metric. Pre-LN
normalization and the AdamW optimizer gave stable optimization, while weight tying supplied
the regularization needed to generalize from a small corpus. The BPE variant scored roughly a
third as well despite reaching a lower training loss.

\section{Error Analysis}

The Pre-LN Transformer succeeded where the Post-LN version had not, because normalizing
before each sub-block kept early training stable instead of diverging at the learning rates
the warmup schedule reaches. The stacked Bi-GRU performed respectably once packing and
padding were optimized, but lacked the capacity to compete.

The dominant error mode in the final model is out-of-vocabulary words. A frequency-2 cutoff
leaves rare words --- proper nouns above all --- outside the vocabulary, where they are
replaced by \texttt{<UNK>} and cannot be recovered. Word dropout during training would force
the model to cope with missing tokens and should improve this; subword modelling would
address it more directly, provided the segmentation is sound.

The most useful insight of the project was that a high-capacity Transformer was the wrong
instrument for this corpus. The larger configurations did not have enough data to learn from
and overfitted rather than generalized. A smaller, well-regularized Transformer outperformed
them --- capacity was never the binding constraint.

\section{Conclusion}

The project's outcome rested on diagnosing architectural instability and controlling
overfitting, not on scaling the model up. My main finding is that for a structurally advanced
architecture like the Pre-LN Transformer, stability and data robustness matter more than raw
capacity. The volatility of the early RNN and Post-LN runs was resolved by adopting the Pre-LN
architecture together with a simple, fixed optimization strategy.

\textbf{Future directions.} A correctly configured subword tokenizer such as BPE should help
appreciably, since Bengali is morphologically rich and word-level vocabularies handle it
badly --- my attempt failed on segmentation quality, not on the premise. Beam search decoding
would likely raise all three metrics at some cost in inference time. Finally, I used a single
architecture for both language pairs throughout; given that Hindi sentences average almost
twice the length of Bengali ones, tuning each pair separately is likely to pay off.

\section*{References}

\begin{enumerate}
    \item Vaswani, A., Shazeer, N., Parmar, N., Uszkoreit, J., Jones, L., Gomez, A. N.,
          Kaiser, {\L}., and Polosukhin, I. (2017). \emph{Attention Is All You Need}. Advances
          in Neural Information Processing Systems 30.
          \url{https://arxiv.org/abs/1706.03762}

    \item Wang, Q., Li, B., Xiao, T., Zhu, J., Li, C., Wong, D. F., and Chao, L. S. (2019).
          \emph{Learning Deep Transformer Models for Machine Translation}. ACL 2019.
          \url{https://arxiv.org/abs/1906.01787}

    \item Press, O., and Wolf, L. (2017). \emph{Using the Output Embedding to Improve Language
          Models}. EACL 2017. \url{https://arxiv.org/abs/1608.05859}

    \item Loshchilov, I., and Hutter, F. (2019). \emph{Decoupled Weight Decay Regularization}.
          ICLR 2019. Establishes that $L_2$ regularization and weight decay are not equivalent
          for adaptive optimizers, which is why AdamW is used here.
          \url{https://arxiv.org/abs/1711.05101}

    \item Hendrycks, D., and Gimpel, K. (2016). \emph{Gaussian Error Linear Units (GELUs)}.
          \url{https://arxiv.org/abs/1606.08415}
\end{enumerate}

\end{document}
